\scanpage{38}
\begin{center}$-\ 27\ -$\end{center}
\hfill n$^\circ$ 371

tion ainsi définie sur $A^S$ est $A$-bilinéaire et vérifie les relations (5) du n$^\circ$ 1 ; en outre, elle est \uline{associative}, car on a
\[
\sum_{uvw=t}\alpha_u\beta_v\gamma_w=\sum_{rw=t}\left(\left(\sum_{uv=r}\alpha_u\beta_v\right)\gamma_w\right)=\sum_{us=t}\left(\alpha_u\left(\sum_{vw=s}\beta_v\gamma_w\right)\right).
\]
Cette multiplication et la structure de $A$-module de $A^S$ \unsure{définissent|définissant} donc sur $A^S$ une structure d'\uline{algèbre associative} par rapport à $A$ ; nous dirons que l'ensemble $A^S$, muni de cette structure, est l'\uline{algèbre large du monoïde} $S$, relative à l'anneau $A$.

Il est immédiat que l'\uline{algèbre} $A^{(S)}$ du monoïde $S$ relative à $A$ (dite encore \uline{algèbre stricte} de $S$ lorsqu'on veut éviter toute confusion) est une \uline{sous-algèbre} de l'algèbre large de $S$ relative à $A$ (et est identique à cette dernière lorsque $S$ est fini). \uline{Par abus de langage}, on note encore tout élément $(\xi_s)_{s\in S}$ de l'algèbre large de $S$ relative à $A$ par la même notation $\sum_{s\in S}\xi_s s$ que les éléments de l'algèbre stricte de $S$ ; bien entendu le signe de sommation qui figure dans cette notation ne représente aucune opération algébrique puisqu'il porte sur une \uline{infinité} de termes $\ne0$. Avec cette notation, la multiplication dans l'algèbre large de $S$ est encore donnée par la formule (23) du n$^\circ$ 11.

Si $S$ est commutatif, il en est de même de son algèbre large $A^S$ ; si $S$ admet un élément neutre $u$, $u$ est aussi élément unité de $A^S$. Si $T$ est une partie stable de $S$, l'algèbre large $A^T$ du monoïde $T$ s'identifie canoniquement à une sous-algèbre de l'algèbre large de $S$.

Parmi les monoïdes satisfaisant à la condition (D), citons en particulier l'ensemble $N$ des entiers naturels, muni de la structure définie par l'\uline{addition}, et l'ensemble $N^*$ des entiers $\ge1$, muni de la structure définie par la \uline{multiplication}. Au chap. IV, nous étudierons l'algèbre stricte (anneau des polynômes sur $A$ à une indéterminée) et l'algèbre
\mspage{37}
